\chapter{Wykład VI}
\section{Skończenie generowane grupy abelowe. Torsja.}

\begin{lemma}
	Niech $\left( A, +\right) $, będzie skończenie generowaną grupą abelową. Niech 
	\begin{align*}
		f: \mathbb{Z}^{n} &\longrightarrow A \\
		(k_1,k_2,\ldots,k_{n}) &\longmapsto f((k_1,k_2,\ldots,k_{n})) = k_1a_1+k_2a_2+\ldots+k_{n}a_{n}
	.\end{align*} 
	
	\begin{enumerate}[label=\arabic*.]
		\item $f$ jest homomorfizmem.
		\item $\im(f) = \langle a_1,a_2,\ldots,a_{n} \rangle \le A$
	\end{enumerate}
\end{lemma} 
\begin{proof}
	TODO
\end{proof} 
\begin{theorem}
	$\exists\left( n\in \N_{>0} \right)\exists\left( H\le \mathbb{Z} ^{n} \right)\left( A\cong \\mathbb{Z}^{n}/H \right)   $
\end{theorem} 
\begin{proof}
	TODO
\end{proof} 
\begin{lemma}
	Niech $B$ jest skonczona i abelowa, $\psi : B \to \mathbb{Z} ^{n}$ jest epimorfizmem.

	Wtedy $B\cong \ker(\psi) \rtimes\mathbb{Z}^{n} $
\end{lemma} 
\begin{proof}
	TODO
\end{proof} 
\begin{theorem}
	$\forall\left( H\le \mathbb{Z} ^{n} \right) \exists\left( 0\le m\le n \right)\left( H\cong\mathbb{Z}^{m} \right)  $
\end{theorem} 
\begin{proof}
	TODO
\end{proof} 
\begin{theorem}
	$H_1\unlhd G_1, H_2\unlhd G_2 \implies H_1\times H_2 \unlhd  G_1 \times G_2 $
	oraz
	 \[
		 \mfaktor{G_1\times G_2}{H_1\times H_2} \cong \mfaktor{G_1}{H_1} \times \mfaktor{G_2}{H_2}
	.\] 	
\end{theorem} 
\begin{proof}
	TODO
\end{proof} 
\begin{eg}
	
	Niech 
	\begin{itemize}
		\item $n=3$, 
		\item $h_1 = \left( 2,0,0 \right) $,
		\item $h_2 = \left( 0,-3,0 \right) $,
		\item $h_3 = \left( 0,0,0 \right) $
	\end{itemize}
	Wtedy $H= 2\mathbb{Z} \times 3\mathbb{Z} \times \{0\} $
Wtedy na mocy twierdzenia
\[
	\mfaktor{\mathbb{Z} ^3}{H} \cong \mfaktor{\mathbb{Z}}{2\mathbb{Z}} \cong \mfaktor{\mathbb{Z} }{\{0\} }
.\] 
Otrzymaliśmy produkt grup cyklicznych.
\end{eg} 

Dla danej macierzy $M\in M_{n\times n}\left( \mathbb{Z}  \right) $ definiujemy \[
\mfaktor{\mathbb{Z} ^{n}}{M} := \mfaktor{\mathbb{Z} ^{n}}{N}, \text{ gdzie }N:= \text{ podgrupa } \mathbb{Z} ^{n}\text{ generowana przez wiersze macierzy M}
.\] 

\begin{lemma}
	Jeśli $M\in M_{n\times n}\left( \mathbb{Z}  \right) $ jest diagonalna to
	\[
	\mfaktor{\mathbb{Z} ^{n}}{M} 
	.\] 
	jest produktem grup cyklicznych.
\end{lemma} 
\begin{proof}
	TODO
\end{proof} 

Szukamy procedury ktora przerabia dowolna macierz M w macierz diagonalna $D$ tak że \[
\mfaktor{\mathbb{Z} ^{n}}{M} \cong \mfaktor{\mathbb{Z} ^{n}}{D} 
.\] 

Jeśli to nam się uda to na mocy lematu(?) dostajemy główną tezę twierdzenia(?)

\begin{lemma}
	Niech $\varphi : G_1 \to G_2 $będzie izomorfizmem grup, $N_1\unlhd G_1$.
	Wtedy $\varphi [N_1] \unlhd G_2$ oraz $\mfaktor{G_1}{N_1} \cong \mfaktor{G_2}{\varphi \left[ N_1 \right]  } $
\end{lemma} 
\begin{proof}
	TODO
\end{proof} 
\begin{lemma}

Weźmy $\varphi \in \aut\left(\mathbb{Z} ^{n}\right)  $. Dla $\begin{bmatrix} m_1 \\ \hdots \\ m_{n} \end{bmatrix} \in M_{n\times n}\left( \mathbb{Z}  \right) $.
Niech $\varphi \left( M \right) := \begin{bmatrix} \varphi (m_1) \\ \vdots \\ \varphi \left( m_{n} \right)    \end{bmatrix} \in M_{n\times n}\left( \mathbb{Z}  \right)  $

Teza:
\[
\mfaktor{\mathbb{Z} ^{n}}{M} \cong \mfaktor{\mathbb{Z} ^{n}}{\varphi \left( M \right)  }.
.\] 
\end{lemma} 
\begin{proof}
	TODO
\end{proof} 	

\begin{lemma}
	Niech  $M\in M_{n\times n}\left( \mathbb{Z}  \right) $ i załóżmy, że $M'$ powstaje z $M$ przez zastosowanie jednej z poniższych operacji:
	\begin{enumerate}[label=\roman*)]
		\item zamiana dwóch wierszy miejscami,
		\item zamiana dwóch kolumn miejscami,
		\item dodanie do wiersza innego wiersza pomnożonego przez liczbę całkowitą,
		\item dodanie do kolumny innej kolumny pomnożonej przez liczbę całkowitą.
	\end{enumerate}
	Wtedy 
	\[
	\mfaktor{\mathbb{Z}^{n}}{M} \cong \mfaktor{\mathbb{Z} ^{n}}{M'} 
	.\] 
\end{lemma} 
\begin{proof}
	TODO
\end{proof} 
\begin{lemma}
	$\forall\left( M\in M_{n\times n}\left( \mathbb{Z}  \right)  \right) \left( \exists\left( D\in M_{n\times n} \left( \mathbb{Z}  \right) \right)  \right) $
\end{lemma} 
\begin{proof}
	TODO (D pozyskujemy przez ciąg operacji z L6)
\end{proof} 

Załóżmy ze $A \cong \mathbb{Z} ^{n}\times \mathbb{Z}_{n_1}\times \mathbb{Z} _{n_2}\times \ldots \times \mathbb{Z}_{n_{k}}$ 

Czy to jest rozkład jednoznaczny?

Nie, bo $\mathbb{Z}_{6} \cong \mathbb{Z}_{2}\times x\mathbb{Z}_{3}$, czyli nie a jednoznaczności.

\begin{theorem}[Bylo jako zadanie]
	NWD( $k$, $l$ )$\iff \mathbb{Z} _{kl}\cong \mathbb{Z}_{k}\times \mathbb{Z}_{l}$
\end{theorem} 
Stąd
\[
	A\cong  \mathbb{Z}^{n}\times \mathbb{Z}_{p_1}^{a_{11}} \times \mathbb{Z}_{p_{1}^{2}}^{a_{12}} \times \ldots \times \text{pogubilem sie w indeksach}
.\] 
gdzie $p_i$ jest liczba pierwszą, $a_{ij}\in \N$


Chcemy pokazać, że istnieje i jest jednoznaczny taki rozkład, z dokładnośćią do kolejnosci czynników.
\begin{theorem}
	Niech $p_1,p_2,\ldots, P_{k}$ będzie ciągiem parami różnch liczb pierwszych. $a_{ij},b_{ij}\in \N$ 
	\[
	\underbrace{\mathbb{Z} _{p_1}^{a_{11}}\times \ldots \times \mathbb{Z} _{p_{1}^{n}}^{a_{1n}}}_{A_1} \times \ldots \times \underbrace{\mathbb{Z} _{p_{k}}^{a_{k_1}}\times \ldots\times \mathbb{Z} _{p_{k}^{n}}^{a_{kn}} }_{A_k} =: A
	,\] 
	oraz
	\[	
		\underbrace{\mathbb{Z} _{p_1}^{b_{11}}\times \ldots \times \mathbb{Z} _{p_{1}^{n}}^{b_{1n}}}_{B_1} \times \ldots \times \underbrace{\mathbb{Z} _{p_{k}}^{b_{k_1}}\times \ldots\times \mathbb{Z} _{p_{k}^{n}}^{b_{kn}} }_{B_{k}} =: B
	.\] 
\end{theorem} 
\begin{proof}
	TODO
\end{proof} 
\begin{definition}[Torsja]
	Niech $A$ będzie grupą abelową
	\[
	\tor(A) := \{a\in A: \text{rząd}\left( a \right) < \aleph_0\} 
	.\] 
\end{definition} 
\begin{remark}
	\begin{enumerate}[label=\arabic*)]
		\item $\tor(A) \le A$
		\item $A$ będąca grupą skończoną $\implies \tor(A) = A$ oraz $\tor\left( \mathbb{Z} ^{n}\times A \right) = \{0\}\times A $
	\end{enumerate}

	Zatem \[
	\mfaktor{\mathbb{Z} ^{n}\times A}{\tor\left( \mathbb{Z} ^{n}\times A \right) } = \mfaktor{\mathbb{Z}^{n}\times A}{\{0\} \times A}\cong \mathbb{Z} ^{n} 
	.\] 
\end{remark} 
\begin{theorem}
	Niech $A,B$ będą grupami skońcoznymi i abelowymi oraz $n,m \in \N$.

	Jeżeli $\mathbb{Z}^{n}\times A \cong \mathbb{Z}^{m} \times B \implies n = m \land A\cong B $
\end{theorem} 
\begin{proof}
	TODO
\end{proof} 
\begin{eg}
	Czy grupy $\mathbb{Z}_{6}\times \mathbb{Z}_{18} \cong \mathbb{Z} _{12}\times \mathbb{Z}_{9}$?
	(Brak notatki)
\end{eg} 
